\section{Sharpness for arbitrary assignment sets}
\label{sec:replacement-sharpness}

For each fixed integer \(k\ge3\), two exponents in
Section~\ref{sec:replacement} are optimal for estimates uniform over
arbitrary assignment sets: the power \(k/(k-1)\) of the inverse killing
probability in the lifetime bound, and the power \(k-1\) of the inverse
block length in the one-clause replacement error.  Both examples below
use subcubes with many frozen coordinates.  A common observation
identifies when a block of clauses can kill such a subcube.

Write \(r=k-1\).  We retain the proper clause law:
a proper \(s\)-clause on \(u\) coordinates chooses \(s\) distinct
coordinates uniformly and gives them independent fair signs.  Different
clauses are independent, with repetition of whole clauses allowed.
Recall that \(q_1(S)\) is the probability that one proper \(r\)-clause
kills \(S\), \(R_t(S)\) is the probability that \(t\) proper \(k\)-clauses
kill \(S\), and \(\tau_k(S)=\EE H_k(S)\).  For nonempty \(S\),
\(H_k(S)\) is the least positive index \(t\) for which the first \(t\)
proper \(k\)-clauses jointly kill \(S\); the convention for the empty set
is \(H_k(\varnothing)=0\).

For integers \(u\ge b\ge k\), consider the nonempty subcube
\[
 S_{u,b}=\{x\in\{0,1\}^u:x_1=\cdots=x_b=0\}.
\]
Its first \(b\) coordinates are frozen and the others are free.  If
\(h_s(S_{u,b})\) denotes the probability that one proper \(s\)-clause
kills this subcube, the frozen-coordinate formula
\eqref{eq:frozen-killing} reads
\begin{equation}\label{eq:sharp-subcube-probabilities}
 h_s(S_{u,b})=\frac{(b)_s}{2^s(u)_s}\quad(1\le s\le u),
 \qquad q_1(S_{u,b})=h_r(S_{u,b}),
\end{equation}
where \((a)_s=a(a-1)\cdots(a-s+1)\).  In particular,
\begin{equation}\label{eq:sharp-successive-ratio}
 \frac{h_k(S_{u,b})}{q_1(S_{u,b})}
 =\frac{b-r}{2(u-r)}.
\end{equation}

For every integer \(t\ge1\), the same subcube satisfies
\begin{equation}\label{eq:sharp-subcube-bounds}
 1-\bigl(1-h_k(S_{u,b})\bigr)^t
 \le R_t(S_{u,b})
 \le t h_k(S_{u,b})+\binom t2\frac{k^2}{u}.
\end{equation}
The lower bound records the independent chances that a single clause
kills the whole subcube.  To prove the upper bound, restrict all \(t\)
clauses to the frozen values.  Each restricted clause is already true,
is the empty false clause, or is a nonempty clause on free coordinates.
If no restricted clause is empty and false, and no two original clauses
share a free coordinate, then the remaining nonempty clauses use
disjoint sets of variables.  Choose one literal in each and set it true,
then assign the other free coordinates arbitrarily.  The resulting
assignment lies in \(S_{u,b}\) and satisfies the block.
Thus a block can kill the subcube only if an individually killing clause
occurs or a pair of clauses shares a free coordinate.  For a fixed
free coordinate, the probability that both independent proper clauses
contain it is \((k/u)^2\).  Summing over the \(u-b\) free coordinates
bounds the probability for a given pair by \(k^2/u\).  A union bound over
clauses and pairs proves the upper bound in
\eqref{eq:sharp-subcube-bounds}.

\begin{proposition}[Optimality of the lifetime exponent]
\label{prop:lifetime-sharp}
Fix an integer \(k\ge3\), and put \(\beta=k/(k-1)\).  For every real
\(\gamma<\beta\) and all constants \(D,A>0\), there exist arbitrarily
large integers \(u\ge k\) and nonempty sets
\(S\subseteq\{0,1\}^u\) such that
\[
 q_1(S)\ge A u^{-(k-1)},\qquad
 \tau_k(S)>Dq_1(S)^{-\gamma}.
\]
Consequently, the killing-probability exponent \(\beta\) in
Proposition~\ref{prop:lifetime} cannot be decreased in a bound uniform
over arbitrary assignment sets,
even above any fixed multiple of the scale \(u^{-(k-1)}\).
\end{proposition}

\begin{proof}
Let the integer \(N\ge2\) tend to infinity, and put
\[
 u=N^{2k+1},\qquad b=N^{2k},\qquad S=S_{u,b}.
\]
Since \(k\) is fixed and \(b/u=N^{-1}\),
\eqref{eq:sharp-subcube-probabilities} gives
\begin{equation}\label{eq:sharp-frozen-probabilities}
 a_N:=q_1(S)\sim 2^{-r}N^{-r},
 \qquad h_N:=h_k(S)\sim 2^{-k}N^{-k}.
\end{equation}
Choose \(t_N=\lfloor1/(4h_N)\rfloor\).  This is positive for all
sufficiently large \(N\), and \(t_N=\Theta_k(N^k)\).  A set survives
the first \(t_N\) clauses exactly when \(H_k(S)>t_N\), so
\eqref{eq:sharp-subcube-bounds} yields
\begin{equation}\label{eq:sharp-survival}
 \PP(H_k(S)>t_N)
 =1-R_{t_N}(S)
 \ge 1-t_Nh_N-\binom{t_N}{2}\frac{k^2}{u}
 \ge\frac34-o(1).
\end{equation}
Here \(t_Nh_N\le1/4\), while \(t_N^2/u=O_k(N^{-1})\).
It follows that
\[
 \tau_k(S)\ge t_N\PP(H_k(S)>t_N)=\Omega_k(N^k).
\]
For the reverse bound, let \(J\) be the index of the first clause that
kills the original set \(S\) by itself.  The cumulative sequence kills \(S\)
no later than \(J\).  Since \(J\) has the geometric distribution with
success probability \(h_N>0\),
\[
 \tau_k(S)\le\EE J=h_N^{-1}=O_k(N^k).
\]
Thus
\[
 \tau_k(S)=\Theta_k(N^k)
          =\Theta_k\bigl(a_N^{-\beta}\bigr).
\]
For each fixed \(\gamma<\beta\), the ratio
\(\tau_k(S)/a_N^{-\gamma}\) tends to infinity, since
\(k-r\gamma>0\).  Finally,
\[
 a_Nu^r\sim 2^{-r}N^{2kr}=2^{-r}b^r\longrightarrow\infty.
\]
For all sufficiently large \(N\), the two displayed requirements
therefore hold for the prescribed \(D\) and \(A\).  The values of \(u\)
are arbitrarily large.
\end{proof}

\begin{proposition}[Optimality of the block-error exponent]
\label{prop:replacement-sharp}
Fix an integer \(k\ge3\).  For integers \(g\ge2\), there are nonempty
sets \(S_g\subseteq\{0,1\}^{u_g}\), with \(u_g\ge k\), for which
\[
 q_1(S_g)-R_g(S_g)\sim 2^{-k}g^{-(k-1)}
 \qquad(g\to\infty).
\]
In particular, no error \(o(g^{-(k-1)})\), uniform over \(u\) and the
assignment set, can replace the error in
Lemma~\ref{lem:one-replacement}.
\end{proposition}

\begin{proof}
For every integer \(g\ge2\), put
\[
 u_g=g^{k+2},\qquad b_g=g^{k+1},\qquad S_g=S_{u_g,b_g},
\]
and abbreviate \(a_g=q_1(S_g)\) and \(h_g=h_k(S_g)\).
Equations~\eqref{eq:sharp-subcube-probabilities} and
\eqref{eq:sharp-successive-ratio} give
\[
 a_g\sim2^{-r}g^{-r},\qquad
 \frac{gh_g}{a_g}=\frac{g(b_g-r)}{2(u_g-r)}
 \longrightarrow\frac12.
\]
The upper bound in \eqref{eq:sharp-subcube-bounds}, and the two-term
inclusion--exclusion lower bound for the independent individual killing
events, give
\[
 gh_g-\binom g2h_g^2
 \le R_g(S_g)
 \le gh_g+\binom g2\frac{k^2}{u_g}.
\]
The collision error is \(O_k(g^{-k})=o(a_g)\).  Also
\(h_g=\Theta_k(g^{-k})\), so the inclusion--exclusion error is
\(O_k(g^{-2(k-1)})=o(a_g)\).  Hence
\[
 R_g(S_g)=gh_g+o(a_g),\qquad
 q_1(S_g)-R_g(S_g)
   =\bigl(\tfrac12+o(1)\bigr)a_g
   \sim2^{-k}g^{-r}.
\]
This is the stated asymptotic.
\end{proof}

These constructions delimit estimates uniform over arbitrary assignment
sets.  They leave open stronger lifetime estimates for classes of sets
with additional structure, and variance arguments using information
beyond the killing probability.  In particular, they do not show that
the logarithm in the three-clause variance upper bound is necessary.
